容易看出空间 $\widetilde L^p(\Omega)$ 不完备 (见 Problem~\ref{ex:continuous-lp-not-complete}). 因此作如下定义.

\begin{Definition}
\label{def:lp-space-completion}
设 $1\leq p<\infty$. 那么 $L^p(\Omega)$ 是 $\widetilde L^p(\Omega)$ 的完备化.
\end{Definition}

\begin{Remark}
\label{rem:lp-completion-defines-distributions}
若序列 $f_n$ 在 $\widetilde L^p(\Omega)$ 中为 Cauchy 序列, 则由 Hölder 不等式, 对每个试验函数 $\phi$, $\int_\Omega f_n\phi\,\mathrm{d}\bm{x}$ 都收敛. 因而 $f_n$ 在分布意义下收敛, 并且由 Theorem~\ref{thm:distribution-sequential-completeness}, $L^p(\Omega)$ 的每个元素都定义了一个分布.
\end{Remark}

需要证明 $L^p(\Omega)$ 的不同元素对应不同的分布.

\begin{Lemma}
\label{lem:lp-completion-faithful-distribution-inclusion}
设 $f_n$ 为 $\widetilde L^p(\Omega)$ 中的 Cauchy 序列 $(1\leq p<\infty)$, 并假设对每个 $\phi\in\mathcal{D}(\Omega)$ 有 $\int_\Omega f_n\phi\,\mathrm{d}\bm{x}\to0$. 那么 $\|f_n\|_p\to0$.
\end{Lemma}

\begin{Proof}
反设 $\lim_{n\to\infty}\|f_n\|_p=M>0$. 选择足够大的 $N$, 使当 $n,m\geq N$ 时 $3M/2>\|f_n\|_p>M/2$ 且 $\|f_n-f_m\|_p<\epsilon$. 这里 $\epsilon$ 取为小量. 由 Problem~\ref{ex:test-functions-dense-lp}, $\mathcal{D}(\Omega)$ 在 $L^p(\Omega)$ 中稠密. 因而可以选择 $\psi\in\mathcal{D}(\Omega)$, 使 $\|f_N-\psi\|_p<\epsilon$; 更具体地, 可以将 $\psi$ 选为多项式与截断函数的乘积. 现在令
\begin{equation}\label{eq:lp-faithful-distribution-dual-test-function}
g=\begin{cases}|\psi|^{p-1}\operatorname{sgn}\psi&p>1\\\operatorname{sgn}\psi&p=1.\end{cases}
\end{equation}
于是 $g$ 是紧支集分片连续函数. 若 $p>1$, 有 $g\in L^q(\Omega)$, 且可以选择 $\phi\in\mathcal{D}(\Omega)$, 使 $\|g-\phi\|_q<\epsilon$; 这里 $q=p/(p-1)$. 现在对 $n\geq N$ 估计如下:
\begin{equation}\label{eq:lp-faithful-distribution-positive-lower-bound}
\begin{aligned}
\left|\int_\Omega f_n\phi\,\mathrm{d}\bm{x}\right|
&\geq\int_\Omega\psi g\,\mathrm{d}\bm{x}\\
&\quad-\left|\int_\Omega\psi(\phi-g)+(f_n-f_N)\phi+(f_N-\psi)\phi\,\mathrm{d}\bm{x}\right|\\
&\geq\int_\Omega|\psi|^p\,\mathrm{d}\bm{x}-\|\psi\|_p\|\phi-g\|_q\\
&\quad+\|f_n-f_N\|_p\|\phi\|_q+\|f_N-\psi\|_p\|\phi\|_q\\
&\geq\|\psi\|_p^p-O(\epsilon)\\
&\geq(M/2-\epsilon)^p-O(\epsilon).
\end{aligned}
\end{equation}
由于对大的 $N$ 可以把 $\epsilon$ 选得任意小, 这与当 $n\to\infty$ 时 \eqref{eq:lp-faithful-distribution-positive-lower-bound} 的左端应趋于零相矛盾.

若 $p=1$, 选取 $\phi$ 使 $\max|\phi|\leq1$, 并使 $\phi$ 除在 $g$ 不连续的那些曲面的一个小邻域内以外, 都一致接近 $g$. 于是可使 $\int\psi(\phi-g)\,\mathrm{d}\bm{x}$ 任意小. 其余与上述论证相同, 只是以初等不等式 $|\int fg\,\mathrm{d}\bm{x}|\leq\max|g|\|f\|_1$ 代替 Hölder 不等式.
\end{Proof}

\begin{Remark}
\label{rem:lp-point-values-and-measure-zero}
考虑序列 $f_n(x)=\exp(-nx^2)$. 容易看出, 当 $1\leq p<\infty$ 时, $f_n\to0$ 在 $\widetilde L^p(\mathbb{R})$ 中成立, 因而该序列可与常值函数 $0\in L^p(\mathbb{R})$ 相认同. 但注意, 对每个 $n$ 都有 $f_n(0)=1$. 这个例子应使读者相信, 不能把 $L^p$ 的元素视为具有点值的函数. 幸运的是, 利用测度论的结果可以证明, ``几乎'' 可以这样做. 在测度论中, $L^p$ 的元素定义为仅在 ``零测集'' 上不同的函数的等价类 (例如, 只在有限个点上不同). 更精确地, $L^p(\Omega)$ 是所有使 $|f|^p$ 可积的 Lebesgue 可测函数 $f$ (的等价类) 的集合. 这一定义与我们的定义等价. 为证明这一等价性, 只需证明有界连续函数在空间 $L^p$ (按积分理论中的定义) 中稠密. 由于未假定读者熟悉 Lebesgue 积分, 这里不再深入讨论这一点.
\end{Remark}

根据对 $L^p(\Omega)$ 的定义, 可以将 ``积分'' 定义为极限. 例如, 若 $f\in L^1(\Omega)$, 则存在序列 $f_n\in\widetilde L^1(\Omega)$, 使 $f_n\to f$, 并可定义
\begin{equation}\label{eq:lp-completion-integral-definition}
\int_\Omega f\,\mathrm{d}\bm{x}=\lim_{n\to\infty}\int_\Omega f_n\,\mathrm{d}\bm{x}.
\end{equation}
此外, 若 $f\in L^p(\Omega)$ 且 $g\in L^q(\Omega)$, 其中 $p,q\in(1,\infty)$ 且 $p^{-1}+q^{-1}=1$, 可以选择序列 $f_n\in\widetilde L^p(\Omega)$ 和 $g_n\in\widetilde L^q(\Omega)$, 使 $f_n\to f$ 且 $g_n\to g$. Hölder 不等式意味着 $f_ng_n$ 在 $L^1(\Omega)$ 中收敛, 自然将极限称为 $fg$. 按此解释, Hölder 不等式对所有 $f\in L^p(\Omega)$ 和 $g\in L^q(\Omega)$ 成立. 这些例子是分布方法的典型做法. 证明 $L^p$ 及类似空间中函数的性质时, 先考察连续 (甚至更光滑的) 函数, 然后考察适当序列的极限.

\Needspace{4\baselineskip}
\begin{Example}
\label{ex:banach-direct-product}
设 $X_1$ 和 $X_2$ 为 Banach 空间. 那么 $X_1\times X_2$ 在范数 $\|(x_1,x_2)\|=\|x_1\|+\|x_2\|$ 下为 Banach 空间. 在这一情形中, 有时将 $X_1$ 与 $X_1\times\{0\}$, $X_2$ 与 $\{0\}\times X_2$ 相认同. 此时用 $x_1+x_2$ 表示 $(x_1,x_2)$.
\end{Example}
